\begin{abstract}
Supervisors of security operations centres must judge periodically whether monitored
organizations detect, investigate, escalate and close security work properly, yet
the usual instruments are self-assessed maturity models and manual record review.
We study whether the operational records an organization submits can instead be
analysed into findings that cite those records, ranked across organizations, decided
by a human supervisor, and bound to that decision cryptographically. We implement
this workflow as SAT-SA, with 16 deterministic rule-based workers, a
saturating risk fusion and a post-quantum signed, hash-chained receipt, and evaluate
it on frozen, hash-verified experiment bundles. On controlled data the mechanisms
behave as designed: every declared finding family was emitted, validation matched
its declared contract in all 245 perturbed conditions, all
13 tested mutations of the decided record were detected, and all
36 injected workflow interruptions recovered. Simple rules are, however,
competitive: fixed-threshold and median-deviation rules tied the fast-closure
detector, and a closure-speed heuristic nearly matched the ranking; stale and
contradictory records were accepted silently. On a public IT incident log, which is
not security-operations data, the risk score was not associated with service-level
misses (Spearman $\rho=\ensuremath{-}0.11$), while slowest median resolution was
($\rho=0.94$); a failure analysis attributes this to construct
mismatch and detector saturation. The evidence supports mechanism-level claims, not
effectiveness in real security operations, for which no human study or expert labels
exist.

\keywords{Security operations \and Supervisory analytics \and Incident management
\and Process conformance \and Human-in-the-loop decision support \and Post-quantum
signatures \and Negative results.}
\end{abstract}
